\chapter{Geometric linear algebra}

\section{On points and vectors}

Linear algebra is an \emph{algebra} of linear things: a line, and an \emph{algebra} is a rule that governs the manipulation and transformation of objects. Hence, linear algebra is a system that allows you to manipulate lines, along with other things that are constructed from lines: planes, cubes, polygons, etc. In some sense, it is a geometry toolkit which allows us to discuss geometry in a more natural way. In this chapter, I'll introduce a few geometrical objects and concepts, including
\begin{enumerate}[noitemsep]
  \item Points and lines,
  \item Vectors, vector spaces, and basis vectors,
  \item Span,
  \item Vectors operations.
\end{enumerate}
This will cover all the prerequisites of linear algebra. Then, we'll discuss what's so \emph{linear} about a
line in \cref{sec:linesanddotproduct}. Linear algebra will take its form in \cref{sec:
  matricesandtransformations}.

\prerequisites{Analytical geometry, algebra: systems of linear equations, trigonometry}

\section{Linear algebra and transformations}

\section{Rotating a graph}

\section{Coordinate dependent transformations}

